\chapter{Structural adaptation and lifecycle governance}
A system that can propose new structures is not thereby entitled to deploy them. Evolutor therefore separates adaptation from authority:
\[
\text{propose}\rightarrow\text{evaluate}\rightarrow\text{canary}\rightarrow\text{promote}
\]
with rejection and rollback as first-class outcomes.

\section{What may adapt}
Candidates can change model weights, adapters, memory policies, retrieval indexes, compiler rules, prompts/policies or topology. These objects have different evaluation and rollback semantics. ``Self-improvement'' is too coarse a systems type.

\section{Immutable artifact lineage}
Every candidate has immutable identity, parent identity, transformation recipe, data/evaluation provenance and compatibility metadata. The lineage forms a DAG. Mutation creates a new artifact; it does not rewrite the active artifact in place.

\section{Lifecycle state machine}
Stages are candidate, evaluated, canary, active, retired and rejected. Promotion is legal only through required gates. The executable registry enforces this state machine.

\section{Evaluation gates}
Let metrics be vector $m$ and required thresholds $g$. A simple gate requires
\[
m_i\ge g_i\quad\forall i,
\]
but real gates include upper bounds, confidence intervals, regression budgets, safety invariants and workload-specific slices. Passing one aggregate score is insufficient.

\section{Continual learning and forgetting}
Sequential adaptation can degrade earlier capabilities. Elastic Weight Consolidation is a concrete method that penalizes movement of parameters judged important to previous tasks:
\[
L(\theta)=L_B(\theta)+\frac{\lambda}{2}\sum_i F_i(\theta_i-\theta_{A,i}^\star)^2.
\]
This is evidence that stability/plasticity can be operationalized; it is not a universal solution. Evolutor requires regression suites regardless of adaptation algorithm.

\section{Candidate creation}
A candidate records parent, transformation, code/data versions, seeds and budgets. Training or editing occurs in an isolated workspace. Active serving artifacts remain immutable.

\section{Shadow and canary}
Shadow evaluation observes production-like inputs without controlling user-visible output. Canary deployment serves a bounded fraction under explicit rollback thresholds. Neither substitutes for offline correctness/safety gates.

\section{Promotion}
Promotion is a durable governance transaction: validate signatures/hashes, gate results, compatibility and policy approval; atomically move the serving pointer; append an audit record. The artifact bytes do not change during promotion.

\section{Rollback}
Rollback selects a previously validated compatible artifact and changes serving state. Runtime state compatibility must be checked: old DOGMA checkpoints or Hermon KV state may be invalid across model/tokenizer changes. Rollback of model version does not imply rollback of irreversible external tool effects.

\section{Structural versus parametric adaptation}
Weight updates are parametric. Adding a memory operator, routing rule or module changes structure. Structural changes require compiler/IR compatibility tests because they can alter execution topology and effect ordering.

\section{Population search}
Multiple candidates may descend from one parent. Selection should compare candidates under matched budgets and held-out evaluation, preserving rejected branches for evidence rather than rewriting history.

\section{Governance authority}
Proposal, evaluation and promotion should be separable authorities. An LLM or evolutionary search may propose a candidate; deterministic policy and signed release mechanisms decide whether it can become active.

\section{Safety invariants}
Some gates are hard constraints rather than objectives: tenant isolation, authorization-before-side-effect, artifact integrity, bounded resources and trace completeness. Optimization cannot trade these away for benchmark score.

\section{Drift and monitoring}
After promotion, monitor workload distribution, quality/safety slices, latency/resource SLOs and anomalous tool behavior. Drift can trigger investigation or rollback, but automated rollback itself must have bounded authority.

\section{Executable registry}
\Listing{reference}
The reference creates immutable lineage IDs, evaluates threshold gates, enforces canary-before-promotion, retires the old active artifact and supports explicit rollback. A candidate that improves quality but fails safety is rejected.

\section{Failure modes}
Failures include in-place mutation, evaluating on training data, one-score promotion, incompatible checkpoint reuse, unsigned artifacts, hidden policy changes, irreversible actions during experiments and deleting failed candidates from the evidence ledger.

\section{Twelve worked exercises}
\begin{enumerate}
\item Why immutable active artifacts? \textbf{Reproducibility and safe rollback.}
\item Candidate parent? \textbf{The exact artifact from which it was derived.}
\item Can a rejected candidate canary? \textbf{No.}
\item Why canary after offline gates? \textbf{It tests bounded production conditions, not basic legality.}
\item EWC addresses what? \textbf{A stability/plasticity aspect of catastrophic forgetting.}
\item Is EWC sufficient governance? \textbf{No.}
\item Structural change example? \textbf{Adding/rerouting a memory or tool operator.}
\item Promotion changes artifact bytes? \textbf{No; it changes the active reference.}
\item Model rollback reverses tool effects? \textbf{No.}
\item Why preserve rejected branches? \textbf{Evidence, diagnosis and reproducibility.}
\item Who may propose versus promote? \textbf{They can be different authorities.}
\item Release rubric? \textbf{Immutable lineage; matched evaluation; hard invariants; canary; signed promotion; compatibility checks; monitoring; rollback; audit trail.}
\end{enumerate}

\section{Bridge}
Chapter 44 begins industrial engineering: packaging, deployment and observable services. Lifecycle governance becomes concrete release artifacts, service contracts, health/readiness, telemetry and deployment automation.
